\documentclass[10pt,a4paper]{amsart}
\usepackage[margin=19mm]{geometry}
\usepackage{amsmath,amssymb,mathtools}
\usepackage[scr=rsfs,cal=euler]{mathalfa}
\usepackage{xfrac}
\usepackage[unicode,hidelinks,bookmarks=false]{hyperref}
\newcommand{\ie}{i.e.\ }
\newcommand{\cf}{cf.\ }
\newcommand{\resp}{respectively\ }
\newcommand{\Subsection}{\subsection}
\providecommand{\nobreakdash}{\nobreak\hbox{-}}
\setlength{\emergencystretch}{2em}
\multlinegap0pt
\usepackage{bm}
\usepackage{bbm}
\usepackage{dsfont}
\usepackage{xspace}
\usepackage{cancel}
\usepackage{mathtools}
\usepackage{amsthm}
\makeatletter
\@ifundefined{lemma}{\newtheorem{lemma}{Lemma}}{}
\makeatother
\newcommand{\br}{\mathbf r}
\newcommand{\charset}{X}
\title[Discrete tensor product BGG sequences: splines and finite el]{Discrete tensor product BGG sequences: splines and finite elements}
\author{Francesca Bonizzoni, Kaibo Hu, Guido Kanschat, Duygu Sap}
\date{Source: arXiv:2302.02434v2 (CC BY 4.0)}
\begin{document}
\maketitle

\noindent\emph{Source and licence.} Discrete tensor product BGG sequences: splines and finite elements. Version arXiv:2302.02434v2 (CC BY 4.0), distributed under the Creative Commons Attribution 4.0 International licence, as recorded in the arXiv licence metadata for this exact version and shown on the abstract page https://arxiv.org/abs/2302.02434v2. Full rights notice: licenses/arxiv-2302.02434-CC-BY-4.0.txt.\par\smallskip

\noindent\textit{Editorial scope.} Counted unit: the Lemma whose source label is \texttt{lem:s} in the article (source0.tex lines 1610--1627, source label \texttt{lem:s}), delivered together with the article's own complete original proof of it (source0.tex lines 1629--1706, one \texttt{proof} environment of 78 lines). No mathematics is written, repaired or re-derived: statement and proof are the article's own text line for line. Required context retained uncounted: def:s (lines 291--295); eqn:form (lines 1605--1608). Every definition, display and label of those lines is retained unchanged. Omitted from this package and not redistributed: 1--290, 296--1604, 1609--1609, 1628--1628, 1707--1778. Nothing inside a retained excerpt is omitted or altered: every retained body is delivered exactly as the article's lines are written, so no omission receipt of any kind accompanies this package. Citations: the retained excerpts carry no citation command at all, so no bibliography entry of the article is transcribed and no reference is redistributed. Reference adaptation: none was needed and none was made; every reference inside the delivered unit resolves inside it, so no unresolved reference or citation is produced. Printed slips kept literal: no printed word, symbol, hypothesis or number of the counted statement or of its proof was altered. Layout and class: the article's own class is \texttt{amsart} with packages including showkeys, amssymb, amsmath, xcolor, tikz, tikz-cd, tikzscale, adjustbox, cite, mathrsfs, mathtools, amsthm, verbatim, subeqnarray; the wrapper is the lane's pinned \texttt{amsart} editorial wrapper at 10pt on A4, which restates the theorem-like environments the retained text needs and adds the article's own macro definitions that the retained text needs (\texttt{\textbackslash br}, \texttt{\textbackslash charset}), with extra wrapper packages (bm, bbm, dsfont, xspace, cancel, mathtools). The delivered numbering is the unit's own. Assets: no retained span contains a figure environment, a graphics-inclusion command, a PDF-page inclusion, a table or a TikZ picture; no image file is copied into this package, the package declares no asset dependency and no image file is redistributed with it. Licence: version arXiv:2302.02434v2 of this article is distributed under the Creative Commons Attribution 4.0 International licence, as recorded in the arXiv licence metadata for this exact version and shown on the abstract page https://arxiv.org/abs/2302.02434v2. A full rights notice is delivered at licenses/arxiv-2302.02434-CC-BY-4.0.txt. Characters: every retained excerpt is pure ASCII, so the delivered engine, XeLaTeX with its default Latin Modern text font, reports no missing character.
\begin{multline}
\label{def:s}
s^{i, j}\mu (v_{0},\cdots,  v_{i})(w_{1}, \cdots, w_{j-1})
\\:=\sum_{l=0}^{i}(-1)^{l} \mu (v_{0},\cdots, \widehat{v_{l}}, \cdots,  v_{i})(v_{l}, w_{1}, \cdots, w_{j-1}), \\ \forall v_{0}, \cdots, v_{i}, w_{1}, \cdots, w_{j-1} \in \mathbb{R}^{n},
\end{multline}

\begin{align}\label{eqn:form}
[dx^{1}\wedge \cdots\wedge& dx^{k}](u_{1}, \cdots, u_{k})\\&\nonumber
=\sum_{l=1}^{k}(-1)^{l+1}dx^{1}(u_{l})[dx^{2}\wedge\cdots\wedge dx^{k}](u_{1}, \cdots, \widehat{u_{l}}, \cdots, u_{k}).
\end{align}

\begin{lemma}\label{lem:s}
Let $\sigma\in\Sigma(k,n)$ and $\tau\in\Sigma(m,n)$. Then
\begin{gather}
    \label{eq:app-lemma-2}
    s^{k,m} \left(dx^\sigma\otimes dx^\tau\right)
    = \sum_{l=1}^m (-1)^{l-1} dx^{\tau_l} \wedge dx^\sigma
    \otimes dx^{\tau_1}\wedge \cdots \wedge \widehat{dx^{\tau_l}}\wedge\cdots\wedge dx^{\tau_m}.
\end{gather}

Let $\bm s\in \charset_k$ and $\bm t\in\charset_m$ be the characteristic vectors of $\sigma$ and $\tau$, respectively. Then, equation~\eqref{eq:app-lemma-2} can be equivalently written as
\begin{multline}
    \label{eq:app-lemma-1}
    s^{k, m}\bigl((dx^{1})^{s_1}\wedge\cdots \wedge (dx^{n})^{s_n}\otimes (dx^{1})^{t_1}\wedge \cdots \wedge (dx^{n})^{t_n}\bigr)\\
    =
    \sum_{l=1}^{n}(-1)^{|\bm{t}|_{l}+1}\delta_{1,t_l}dx^{l}\wedge (dx^1)^{s_1}\wedge\cdots \wedge (dx^n)^{s_n}\otimes \cdots\\
    (dx^1)^{t_1}\wedge \cdots \wedge \widehat{dx^l}\wedge\cdots\wedge (dx^n)^{t_n}.
\end{multline}
\end{lemma}

\begin{proof}
Let $\sigma\in\Sigma(k,n)$ and let $\tau\in\Sigma(m,n)$.
By the definition of $s^{k,m}$ in equation~\eqref{def:s}, there holds
\begin{align*}
T_1 &:=
    s^{k,m}\left[dx^\sigma\otimes dx^\tau\right]
    (e_{\br_1},\dots,e_{\br_{k+1}})
    (e_{\beta_1},\dots,e_{\beta_{m-1}})\\
    &=\sum_{\ell=1}^{k+1} (-1)^{\ell+1} \left[dx^\sigma\otimes dx^\tau\right]
    (e_{\br_1},\dots,\widehat{e_{\br_\ell}},\dots, e_{\br_{k+1}})
    (e_{\br_\ell},e_{\beta_1},\dots,e_{\beta_{m-1}})\\
    &=\sum_{\ell=1}^{k+1} (-1)^{\ell+1} 
    dx^\sigma (e_{\br_1},\dots,\widehat{e_{\br_\ell}},\dots, e_{\br_{k+1}})
    dx^\tau (e_{\br_\ell},e_{\beta_1},\dots,e_{\beta_{m-1}}),
\end{align*}
for $\br\in\Sigma(k+1,n)$ and $\beta\in\Sigma(m-1,n)$. We observe that at most one term in the sum on the right is nonzero, namely where $\ell$ can be chosen such there holds
\begin{gather}
    \{\br_1,\dots,\br_{k+1}\}
    = \{\br_\ell,\sigma_1,\dots,\sigma_k\},
    \qquad
    \{\br_\ell,\beta_1,\dots,\beta_{m-1}\}
    =\{\tau_1,\dots,\tau_m\}.
\end{gather}
In particular, we note that $\br_\ell\neq\sigma_i$
for $i=1,\dots,k$ and there is an index $j$ such that $\br_\ell = \tau_j$. Define $\tau'\in\Sigma(m-1,n)$ by removing $\tau_j$ from $\tau$. Then, by~\eqref{eqn:form},
\begin{align*}
    dx^\tau(e_{\br_\ell},e_{\beta_1},\dots,e_{\beta_{m-1}})
    &=
    (-1)^{j+1} [dx^{\tau_j}\wedge dx^{\tau'}](e_{\br_\ell},e_{\beta_1},\dots,e_{\beta_{m-1}})\\
    &= (-1)^{j+1} dx^{\tau_j}(e_{\br_\ell})
    dx^{\tau'}(e_{\beta_1},\dots,e_{\beta_{m-1}})\\
    &= (-1)^{j+1}
    dx^{\tau'}(e_{\beta_1},\dots,e_{\beta_{m-1}}).
\end{align*}
Thus,
\begin{gather}
    \label{eq:app-T1-final}
    T_1 = (-1)^{\ell+j} \delta_{\br_\ell,\tau_j}
    dx^\sigma(e_{\br_1},\dots,\widehat{e_{\br_\ell}},\dots, e_{\br_{k+1}})
    dx^{\tau'}(e_{\beta_1},\dots,e_{\beta_{m-1}}).
\end{gather}
Next, we start from~\eqref{eq:app-lemma-2} and let (again using $\tau'$ for $\tau$ without $\tau_j$)
\begin{align*}
    T_2 &:= 
    \sum_{j=1}^m (-1)^{j+1}
    \left[dx^{\tau_j} \wedge dx^\sigma \otimes dx^{\tau'}\right]
    (e_{\br_1},\dots,e_{\br_{k+1}})
    (e_{\beta_1},\dots,e_{\beta_{m-1}})\\
    &= \sum_{j=1}^m (-1)^{j+1}
    \left[dx^{\tau_j} \wedge dx^\sigma\right]
    (e_{\br_1},\dots,e_{\br_{k+1}})
    dx^{\tau'}
    (e_{\beta_1},\dots,e_{\beta_{m-1}}).
\end{align*}
The last sum contains at most one nonzero term, namely if $j$ can be chosen such that there holds
\begin{gather}
    \{\tau_j,\beta_1,\dots,\beta_{m-1}\}
    =\{\tau_1,\dots,\tau_m\},
    \qquad
    \{\br_1,\dots,\br_{k+1}\}
    =\{\tau_j,\sigma_1,\dots,\sigma_k\}.
\end{gather}
Since $\br_1,\dots,\br_{k+1}$ is as in~\eqref{eq:app-T1-final}, the index not contained in $\sigma$ is $\br_\ell$. By~\eqref{eqn:form} we get
\begin{gather}
    [dx^{\tau_j} \wedge dx^\sigma]
    (e_{\br_1},\dots,e_{\br_{k+1}})
    = (-1)^{\ell+1} dx^{\tau_j}(e_{\br_\ell})
    dx^\sigma(e_{\br_1},\dots,\widehat{e_{\br_\ell}},\dots,e_{\br_{k+1}}).
\end{gather}
Summarizing, we obtain
\begin{gather}
    \label{eq:app-T2-final}
    T_2 = (-1)^{\ell+j} \delta_{\br_\ell,\tau_j}
    dx^\sigma(e_{\br_1},\dots,\widehat{e_{\br_\ell}},\dots, e_{\br_{k+1}})
    dx^{\tau'}(e_{\beta_1},\dots,e_{\beta_{m-1}}),
\end{gather}
which is equal to $T_1$ in~\eqref{eq:app-T1-final}.
\end{proof}

\end{document}
